
Finite electro-thermo-mechanical constitutive models are subject to two independent sets of restrictions: those arising from thermodynamics and those required for mathematical well-posedness. While existing formulations generally satisfy the First and Second Laws of Thermodynamics, recent work by Bonet and Gil \cite{BonetGil2026} has shown that constitutive models based on constant, or purely temperature-dependent, specific heat are fundamentally incompatible with the Third Law of Thermodynamics. Conversely, thermodynamic admissibility alone does not ensure constitutive stability or robust numerical performance.

This paper establishes a polyconvex electro-thermo-visco-hyperelastic constitutive framework for electro-active polymers that simultaneously satisfies the three laws of thermodynamics and the mathematical requirements associated with finite-strain elasticity. The framework makes three principal contributions. \emph{First}, it introduces a class of electro-thermo-hyperelastic constitutive models that satisfy the First, Second and Third Laws of Thermodynamics while preserving polyconvexity for arbitrary temperature, ensuring a strictly positive heat capacity and stable wave propagation. \emph{Second}, it extends this constitutive structure to finite-strain viscoelasticity through a thermodynamically consistent equilibrium/non-equilibrium decomposition that naturally accommodates anisotropic extensions via structural tensors. \emph{Third}, it provides a comprehensive calibration against available thermo-mechanical, electro-mechanical and coupled electro-thermo-mechanical experiments, followed by large-scale finite element simulations that quantify the influence of viscous dissipation and internally generated heat on the coupled response. The resulting framework provides a unified constitutive basis for the predictive simulation of electro-active polymers undergoing finite deformations under coupled electro-thermo-mechanical loading
